
这样，我们通过分叉论证得到了
$A_{n,K}$ 同态 (7.2.363)、
其 $A_\infty$ 版本及相应的同伦。
与配边论证相比，分叉论证有以下三点不足：

\par\noindent\textup{(7.2.364.1)}
如上所述，我们还需要多迭代一层同伦。
也就是说，配边论证只需做到同伦的同伦，
而分叉论证还需要同伦的同伦的同伦。
因此，写出证明细节会稍显繁琐，
尽管我们已有足够的工具完成这一步。

\par\noindent\textup{(7.2.364.2)}
在第 5.3.2 小节的情形中，
仅用分叉论证构造滤过 $A_\infty$ 双模同态
存在一些困难。本小节稍后将详述这一点。

\par\noindent\textup{(7.2.364.3)}
在本书写作时，作者尚无法将分叉论证用于
第 7.5 节给出的定理 X 的证明。

另一方面，与分叉论证相比，
配边论证也有以下不足：

\par\noindent\textup{(7.2.365.1)}
研究第 3.8 节引入的算子
$\mathfrak q$ 和 $\mathfrak r$ 时，
分叉论证能给出更好的结果。
我们将在第 7.4 节说明这一点。

\par\noindent\textup{(7.2.365.2)}
分叉论证只需使用“逐时积”（\emph{twp}）型模空间，
完全不必使用“时间有序积”（\emph{top}）模空间。
因此，采用分叉论证时，部分论证较为简单。

\par\noindent\textup{(7.2.365.3)}
无需完成第 7.3.10 小节给出的
$\operatorname{Incl}$ 几何构造。
这一构造相当精细，因而若能避开，自是可取。

上述差异都偏于技术层面。
只要多费些功夫，很可能一种方法所能证明的结果，
另一种方法也能证明；其中 (7.2.364.2)
似乎是最严重的一点。
不过，在写作本书时，由于
(7.2.364.2)、(7.2.364.3) 和
(7.2.365.1) 的问题，
作者仍需同时使用两种方法，
才能得到尽可能好的结果。
整个构造又相当精细，
因此有几种不同的证明也更为妥当。

下面证明这两种方法得到的
$A_\infty$ 同态相互同伦。
考虑定理~\ref{thm:4-6-1} 的情形，
并沿用其中的记号。

\begin{proposition}{7.2.366}\label{prop:7-2-366}
第~\ref{subsec:7-2-9}、\ref{subsec:7-2-11}
小节得到的滤过 $A_\infty$ 同态
$\psi_*\colon C(L;\Lambda_{0,\mathrm{nov}})
\to C(L';\Lambda_{0,\mathrm{nov}})$，
与本小节得到的同态同伦。
\end{proposition}

\noindent\textbf{证明。}
我们只证明命题的 $A_{n,K}$ 版本。
先取常值族 $\equiv J$，得到图式~\ref{diag:7-2-36}：

\par\medskip\noindent\begin{minipage}{\textwidth}
\centering
$C(L,\lambda)
\xleftarrow{\ \operatorname{Eval}_{\rho=0}\ }
C([0,1]\times L;\lambda)
\xrightarrow{\ \operatorname{Eval}_{\rho=1}\ }
C(L,\lambda)$
\par\smallskip
{\renewcommand{\thefigure}{7.2.36}\refstepcounter{figure}
\normalfont\bfseries 图式~\thefigure\label{diag:7-2-36}}
\end{minipage}\par\medskip

再考虑殆复结构的双参数族
\[
\{J_{\sigma,\rho}\}_{0\le\sigma\le1,\,0\le\rho\le1}
=\{J_{\sigma,\rho}\}_{\sigma,\rho},
\]
它满足 $J_{0,\rho}=\psi_*J$、
$J_{1,\rho}=J_\rho$ 和
$J_{\sigma,0}=\psi_*J$。
利用第~\ref{subsec:7-2-10} 小节定义的模空间
\[
\mathcal M^{\mathrm{main}}_{k+1}
(\{J_{\sigma,\rho}\}_{\sigma,\rho}:\beta;
\operatorname{top}(\sigma),\operatorname{twp}(\rho)),
\]
便得到下列同伦交换图式的纵向箭头：

\par\medskip\noindent\begin{minipage}{\textwidth}
\centering
\includegraphics[width=.93\textwidth]{chapters/ch07-transversality/figures/phys-581-diagram-7-2-37-math-clean-crop.png}
\par\smallskip
{\renewcommand{\thefigure}{7.2.37}\refstepcounter{figure}
\normalfont\bfseries 图式~\thefigure\label{diag:7-2-37}}
\end{minipage}\par\medskip

由构造可知，第三列给出
第~\ref{subsec:7-2-9}、\ref{subsec:7-2-11}
小节得到的同态 $\psi_*$；
第二行导出本小节得到的同态 $\psi_*$。
第一列和第一行都导出恒等映射。
命题得证。\qed
